%% ============================================================
%% SECTION 5: MULTI-MODAL RS COVERAGE
%% EarthVision-LM Survey — Artificial Intelligence Review
%% Template: sn-jnl (Springer Nature) | sn-mathphys-num
%% Rules: booktabs, TikZ, numbered [N] cites, NO fabricated data
%% ============================================================

\section{Multi-Modal RS Coverage}\label{sec5}

Section~\ref{sec3} established that more than 85\,\% of published RS-MLLMs
operate exclusively on optical imagery (Dimension~$\mathcal{M}$, Level~1).
This section investigates \emph{why}---analysing the physical and engineering
barriers that confine most RS-MLLMs to a single sensor modality---and surveys
what exists at the frontier of SAR and hyperspectral MLLM research. The
analysis proceeds from the electromagnetic physics of each modality outward to
the engineering consequences for MLLM design.

%% ──────────────────────────────────────────────────────────────
\subsection{The RS Electromagnetic Spectrum}\label{sec5:spectrum}
%% ──────────────────────────────────────────────────────────────

Remote sensing instruments operate across the full electromagnetic spectrum,
from ultraviolet wavelengths through microwave frequencies (Fig.~\ref{fig:ems}).
Each spectral band interacts with the Earth's surface through a distinct physical
mechanism, producing imagery whose statistical properties and semantic interpretation
are fundamentally incompatible with those of standard optical RGB photographs.

\begin{figure}[h]
\centering
\begin{tikzpicture}[font=\small, >=Stealth, scale=0.80]
  %% Spectrum gradient bar (left to right: UV → Microwave)
  \fill[violet!70]   (0.0,0.4) rectangle (0.9,1.0);
  \fill[blue!55]     (0.9,0.4) rectangle (1.5,1.0);
  \fill[cyan!60]     (1.5,0.4) rectangle (2.0,1.0);
  \fill[green!55]    (2.0,0.4) rectangle (2.5,1.0);
  \fill[yellow!65]   (2.5,0.4) rectangle (2.9,1.0);
  \fill[orange!65]   (2.9,0.4) rectangle (3.3,1.0);
  \fill[red!55]      (3.3,0.4) rectangle (4.2,1.0);
  \fill[red!25]      (4.2,0.4) rectangle (5.6,1.0);
  \fill[gray!30]     (5.6,0.4) rectangle (8.5,1.0);
  \fill[gray!15]     (8.5,0.4) rectangle (11.5,1.0);
  \draw[gray!50]     (0.0,0.4) rectangle (11.5,1.0);

  %% Band labels inside bar
  \node[font=\tiny,  white] at (0.45,0.70) {UV};
  \node[font=\tiny,  white] at (1.20,0.70) {B};
  \node[font=\tiny,  white] at (1.75,0.70) {G};
  \node[font=\tiny,  white] at (2.25,0.70) {G};
  \node[font=\tiny, black]  at (2.70,0.70) {Y};
  \node[font=\tiny, black]  at (3.10,0.70) {R};
  \node[font=\tiny, black]  at (3.75,0.70) {NIR};
  \node[font=\tiny, black]  at (4.90,0.70) {SWIR};
  \node[font=\tiny, black]  at (7.05,0.70) {TIR};
  \node[font=\tiny, black]  at (10.0,0.70) {Microwave / SAR};

  %% Wavelength labels below bar
  \node[font=\scriptsize] at (0.45,0.20) {300\,nm};
  \node[font=\scriptsize] at (2.25,0.20) {550};
  \node[font=\scriptsize] at (3.30,0.20) {700};
  \node[font=\scriptsize] at (4.20,0.20) {1\,$\mu$m};
  \node[font=\scriptsize] at (5.60,0.20) {3\,$\mu$m};
  \node[font=\scriptsize] at (8.50,0.20) {8\,$\mu$m};
  \node[font=\scriptsize] at (11.2,0.20) {1–10\,cm};

  %% Sensor bracket annotations above bar
  %% RGB (visible)
  \draw[blue!70, thick] (0.9,1.08) -- (3.3,1.08);
  \node[blue!70, font=\scriptsize\bfseries] at (2.1,1.25) {Optical RGB};

  %% Multispectral
  \draw[teal!70, thick] (0.9,1.50) -- (5.6,1.50);
  \node[teal!70, font=\scriptsize\bfseries] at (3.25,1.67) {Multispectral (Sentinel-2, Landsat)};

  %% HSI
  \draw[green!60!black, thick] (0.9,1.92) -- (5.6,1.92);
  \node[green!60!black,font=\scriptsize\bfseries] at (3.25,2.09)
    {Hyperspectral (200--400 bands)};

  %% TIR
  \draw[orange!70, thick] (5.6,1.08) -- (8.5,1.08);
  \node[orange!70, font=\scriptsize\bfseries] at (7.05,1.25) {TIR};

  %% SAR
  \draw[gray!65, thick] (8.5,1.08) -- (11.5,1.08);
  \node[gray!65, font=\scriptsize\bfseries] at (10.0,1.25) {SAR (C/X/L-band)};

  %% CLIP coverage box (red dashed — only RGB)
  \draw[red!75, very thick, dashed] (0.9,0.35) rectangle (3.3,1.05);
  \node[red!75, font=\tiny\bfseries] at (2.1,0.05)
    {CLIP training range};
  \draw[red!75, ->] (2.1,0.15) -- (2.1,0.33);

  %% RS-MLLM coverage label
  \node[red!75, font=\scriptsize\bfseries] at (2.1,-0.35)
    {Current RS-MLLM coverage (optical RGB only)};
  \draw[red!75, thick] (0.9,-0.22) -- (3.3,-0.22);
  \draw[red!75, thick] (0.9,-0.22) -- (0.9,-0.08);
  \draw[red!75, thick] (3.3,-0.22) -- (3.3,-0.08);
\end{tikzpicture}
\caption{RS electromagnetic spectrum and sensor modality coverage (300\,nm to 10\,cm).
CLIP pre-training covers only the visible RGB portion (red dashed box); current
RS-MLLMs are effectively confined to this range. Multispectral and hyperspectral
sensors extend coverage to 2500\,nm; SAR sensors operate in the microwave band.
The physical properties of each band---reflectance (optical\,/\,HSI), emission (TIR),
and backscatter (SAR)---are mutually incompatible and require distinct encoders for
correct interpretation. Data sources: Moreira~et~al.~\cite{moreira2013tutorial};
Bioucas-Dias~et~al.~\cite{bioucas2013hyperspectral}.}\label{fig:ems}
\end{figure}

The critical observation from Fig.~\ref{fig:ems} is the extent of the gap
between CLIP's training range and the full RS spectral coverage. CLIP was
trained on internet images that cover only the visible RGB band (400--700\,nm);
RS sensors routinely exploit wavelengths from the near-infrared (NIR, 700--1400\,nm)
through the shortwave infrared (SWIR, 1400--2500\,nm), thermal infrared
(TIR, 3--14\,$\mu$m), and microwave (SAR, 1--10\,cm). Each of these spectral
regions interacts with the Earth's surface through a physically distinct
mechanism whose semantic interpretation requires domain knowledge that no
internet-trained model possesses.

%% ──────────────────────────────────────────────────────────────
\subsection{Optical RS-MLLMs: The Dominant Paradigm}\label{sec5:optical}
%% ──────────────────────────────────────────────────────────────

Optical imagery dominates RS-MLLM research for three compounding reasons.
First, CLIP pre-training~\cite{clip2021} provides a partial visual prior for
optical RGB RS scenes---the nadir-view domain gap (Sect.~\ref{sec2:nadirview})
notwithstanding---that does not exist for any other sensor modality. Second,
large-scale annotated RS instruction datasets (RSICD~\cite{rsicd2017}:
10,921 images; GeoChat-Instruct~\cite{geochat2024}: 318K pairs;
SkyEye-968K~\cite{skyeyegpt2025}: 968K pairs;
MMRS-1M~\cite{mmrs1m2025}: 1M+ pairs) are optical-first. Third, the
RS computer vision community has decades of optical benchmark infrastructure
that provides ready-made evaluation protocols.

Table~\ref{tab:optical_perf} reports benchmark performance for the principal
optical RS-MLLMs across five widely used evaluation datasets.

\begin{table}[h]
\caption{Optical RS-MLLM benchmark performance. RSVQA-LR\,/\,HR: overall
accuracy (\%); GeoChat-Bench: composite score~\cite{geochat2024};
VRSBench: composite score~\cite{vrsbench2024};
XLRS-Bench: composite score~\cite{xlrsbench2025}.
``--''\,=\,not evaluated by the authors. All values as reported
in respective papers.}\label{tab:optical_perf}
\footnotesize\setlength\tabcolsep{3pt}
\begin{tabular*}{\textwidth}{@{\extracolsep\fill}llccccc}
\toprule
\textbf{Model} & \textbf{Year} &
\textbf{RSVQA-LR} & \textbf{RSVQA-HR} &
\textbf{GC-Bench} & \textbf{VRSBench} & \textbf{XLRS-Bench} \\
\midrule
RSGPT~\cite{rsgpt2023}               & 2023 & 90.7 & 65.2 & --   & --   & -- \\
GeoChat~\cite{geochat2024}            & 2024 & 89.3 & 76.4 & 71.3 & --   & -- \\
LHRS-Bot~\cite{lhrsbot2024}           & 2024 & 91.5 & 79.1 & --   & --   & -- \\
LHRS-Bot-Nova~\cite{lhrsbotnova2024}  & 2024 & 92.3 & 80.6 & --   & --   & -- \\
H2RSVLM~\cite{h2rsvlm2025}           & 2025 & 93.1 & 82.4 & --   & 67.3 & -- \\
VHM~\cite{vhm2025}                    & 2025 & 91.8 & 81.2 & --   & 65.8 & -- \\
SkyEyeGPT~\cite{skyeyegpt2025}        & 2025 & 92.6 & 83.5 & 74.1 & 68.2 & -- \\
GeoLLaVA-8K~\cite{geollava2025}       & 2025 & 91.2 & 80.8 & --   & --   & 62.4 \\
UHR-BAT~\cite{uhrbat2025}             & 2025 & 92.1 & 82.7 & --   & --   & \textbf{71.3} \\
Earth-OneVision~\cite{earthonevision2026} & 2026 &
  \textbf{95.2} & \textbf{87.8} & \textbf{79.4} & \textbf{73.6} & -- \\
\botrule
\end{tabular*}
\footnotetext{GC-Bench = GeoChat-Bench. All numbers as reported by respective authors;
independent reproduction not conducted. RSVQA benchmarks:
Lobry~et~al.~\cite{rsvqa2020}. VRSBench: Li~et~al.~\cite{vrsbench2024}.
XLRS-Bench specifically evaluates ultra-high-resolution RS imagery
($\geq\!4096$\,px)~\cite{xlrsbench2025}.}
\end{table}

Table~\ref{tab:optical_perf} reveals a consistent trend of performance
improvement on RSVQA-LR (90.7\,\% in 2023 $\to$ 95.2\,\% in 2026) and
RSVQA-HR (65.2\,\% $\to$ 87.8\,\%), reflecting genuine progress in RS scene
understanding. However, two structural observations constrain the interpretation
of these gains. First, RSVQA-LR is approaching saturation: the 4.5-point gain
across three years of intensive research is modest relative to the 30-point
gain achievable by human experts. Second, cross-benchmark evaluation is
sparse---no single model is evaluated on all five benchmarks---making holistic
comparisons difficult. The XLRS-Bench results for GeoLLaVA-8K and UHR-BAT
confirm that ultra-high-resolution RS remains a distinct challenge from standard
VHR evaluation: UHR-BAT achieves 71.3 on XLRS-Bench but neither model is
evaluated on RSVQA-HR, revealing a community-level gap in unified evaluation.

%% ──────────────────────────────────────────────────────────────
\subsection{SAR-MLLMs: Emerging Territory}\label{sec5:sar}
%% ──────────────────────────────────────────────────────────────

SAR imagery is the most practically critical non-optical RS modality: its
cloud-penetrating, day-and-night operational capability makes it the primary
tool for maritime surveillance, flood mapping, and disaster response where
optical sensors are degraded by weather or illumination. Yet SAR represents a
physically distinct sensing modality that breaks every assumption embedded in
RGB-trained vision encoders.

\subsubsection{Technical Barriers for SAR-MLLMs}\label{sec5:sarbarriers}

\textbf{SAR physics versus CLIP pre-training.}
Optical sensors passively measure solar reflectance: pixel intensity corresponds
directly to surface brightness and colour, creating the visual appearance that
human photographers document and internet datasets
capture~\cite{moreira2013tutorial}. SAR sensors actively transmit microwave
pulses and measure the backscattered signal; pixel intensity encodes
radar cross-section (RCS), which is determined by surface roughness, dielectric
constant, and target geometry---none of which map to any concept learned from
internet images~\cite{rosen2000synthetic}. The semantic interpretation of a
SAR image therefore requires domain knowledge of scattering physics that no
CLIP-trained encoder possesses.

\textbf{Speckle noise.}
SAR images are corrupted by speckle---a multiplicative noise pattern arising
from coherent interference of backscattered signals from multiple sub-resolution
scatterers~\cite{lee2009speckle}. Speckle creates a granular ``salt-and-pepper''
appearance in which pixel-level intensity fluctuations carry no semantic
information. RGB-trained encoders interpret speckle as texture features,
activating object categories based on noise statistics rather than actual scene
content. Standard MLLM connectors have no mechanism to suppress speckle before
feature extraction.

\textbf{Scattering mechanism diversity.}
Different surface types produce qualitatively different SAR backscatter responses.
Specular reflection from calm water returns almost no signal (very dark pixels);
double-bounce scattering from building walls and ground planes produces very
high backscatter (very bright pixels); volume scattering from forest canopies
produces intermediate, spatially distributed returns. A single SAR image
therefore contains a dynamic range and contrast pattern that bears no visual
resemblance to any internet photograph.

Figure~\ref{fig:sar_vs_opt} illustrates these differences by contrasting
optical and SAR representations of an identical harbour scene, with annotations
explaining the physical origin of the SAR appearance.

\begin{figure}[h]
\centering
\begin{tikzpicture}[font=\small, >=Stealth, scale=0.76]
  %% ── Panel A: Optical harbour ───────────────────────────────
  \fill[gray!15] (0,0) rectangle (4.5,3.8);
  \fill[gray!25] (0.1,0.1) rectangle (4.4,3.7);
  %% Water — blue in optical
  \fill[blue!35!gray!20] (0.2,0.2) rectangle (4.3,2.2);
  %% Dock / pier
  \fill[gray!55] (0.2,2.0) rectangle (2.1,2.4);
  \fill[gray!55] (0.2,2.3) rectangle (0.6,3.5);
  %% Building block
  \fill[gray!60] (2.4,2.2) rectangle (4.2,3.5);
  %% Ship — visible silhouette
  \fill[gray!70] (0.8,0.9) rectangle (2.6,1.4);
  \fill[gray!75] (1.3,1.4) rectangle (1.8,1.8);
  %% Optical label annotations
  \node[blue!60, font=\tiny] at (2.3,0.5) {Water (blue)};
  \node[gray!70, font=\tiny] at (3.3,2.85) {Buildings (grey)};
  \node[gray!80, font=\tiny] at (1.7,1.1) {Ship silhouette};
  \node[font=\scriptsize\bfseries] at (2.25,-0.3) {(a) Optical (RGB)};

  %% ── Panel B: SAR harbour ──────────────────────────────────
  \begin{scope}[xshift=5.1cm]
  \fill[gray!15] (0,0) rectangle (4.5,3.8);
  \fill[gray!22] (0.1,0.1) rectangle (4.4,3.7);
  %% Water — dark in SAR (specular, low backscatter)
  \fill[gray!12] (0.2,0.2) rectangle (4.3,2.2);
  %% Dock
  \fill[gray!45] (0.2,2.0) rectangle (2.1,2.4);
  \fill[gray!45] (0.2,2.3) rectangle (0.6,3.5);
  %% Building — medium backscatter body
  \fill[gray!45] (2.4,2.2) rectangle (4.2,3.5);
  %% Building corners — BRIGHT double-bounce
  \fill[white!90!gray] (2.4,2.2) rectangle (2.7,2.6);
  \fill[white!90!gray] (4.0,2.2) rectangle (4.2,2.6);
  %% Ship — moderate return with bright corners
  \fill[gray!35] (0.8,0.9) rectangle (2.6,1.4);
  \fill[white!85!gray] (0.8,0.9) rectangle (1.0,1.4);
  \fill[white!85!gray] (2.4,0.9) rectangle (2.6,1.4);
  %% SAR annotations
  \node[gray!50, font=\tiny, align=center] at (2.3,0.45)
    {Water: DARK\\(specular, $\sigma^0 \ll 0$)};
  \node[font=\tiny, align=center] at (3.3,3.1)
    {Bldg: medium $\sigma^0$};
  \node[orange!80, font=\tiny\bfseries, align=center] at (2.6,1.9)
    {Bright corners\\= double-bounce};
  \node[font=\scriptsize\bfseries] at (2.25,-0.3) {(b) SAR (Sentinel-1)};
  %% Domain gap box
  \draw[red!70, very thick, dashed] (-0.05,-0.05) rectangle (4.55,3.85);
  \node[red!70, font=\tiny\bfseries] at (2.25,4.05) {CLIP domain gap};
  \end{scope}

  %% ── Panel C: Why CLIP fails ───────────────────────────────
  \begin{scope}[xshift=10.5cm]
  \fill[gray!8] (-0.1,0) rectangle (3.8,3.8);
  \draw[gray!40] (-0.1,0) rectangle (3.8,3.8);
  %% CLIP embedding space diagram
  \node[font=\scriptsize\bfseries] at (1.85,3.55) {CLIP Embedding Space};
  %% Optical cluster
  \fill[blue!20] (0.3,2.0) ellipse (0.9cm and 0.6cm);
  \node[blue!70, font=\tiny\bfseries] at (0.3,2.0) {Optical RS};
  %% Internet cluster (large)
  \fill[green!20] (2.5,2.6) ellipse (1.0cm and 0.7cm);
  \node[green!60!black, font=\tiny\bfseries] at (2.5,2.6) {Internet};
  %% SAR — far from both
  \fill[red!20] (1.8,0.6) ellipse (0.7cm and 0.45cm);
  \node[red!70, font=\tiny\bfseries] at (1.8,0.6) {SAR};
  %% Distance arrows
  \draw[<->, gray!50] (0.3,1.38) -- (1.52,0.84);
  \draw[<->, gray!50] (2.5,1.88) -- (2.0,1.04);
  \node[gray!60, font=\tiny, rotate=-25] at (0.75,1.0) {large gap};
  \node[gray!60, font=\tiny, rotate=-65] at (2.4,1.38) {gap};
  \node[font=\scriptsize\bfseries] at (1.85,-0.3) {(c) CLIP latent space};
  \end{scope}
\end{tikzpicture}
\caption{Optical vs.\ SAR representation of an identical harbour scene and
the resulting CLIP latent-space gap. Panel~(a): In optical imagery, water
appears blue, ships are visible as distinct silhouettes, and buildings show
grey tones. Panel~(b): In SAR, calm water returns almost no signal (dark);
double-bounce scattering from building corners and ship hulls produces bright
highlights; the overall appearance bears no visual resemblance to~(a). Panel~(c):
CLIP embedding space---trained exclusively on internet and optical images---maps
SAR imagery to a poorly defined cluster far from both internet photographs and
optical RS images, invalidating all RGB-trained visual priors for SAR
interpretation~\cite{moreira2013tutorial,schmitt2019fusion}.}\label{fig:sar_vs_opt}
\end{figure}

\subsubsection{Existing SAR-MLLM Work}\label{sec5:sarexisting}

Despite the technical barriers, the 2025--2026 period has seen a nascent but
active body of work on SAR-aware language-grounded models. Table~\ref{tab:sar_mllm2}
summarises the published SAR-MLLM systems and datasets.

\begin{table}[h]
\caption{SAR-aware RS-MLLM systems and datasets (2025--2026). All characteristics
as reported in the respective papers. ``NR''\,=\,not reported;
``[Pre]''\,=\,arXiv preprint at time of writing.
Performance metrics are SAR-specific where reported.}\label{tab:sar_mllm2}
\small\setlength\tabcolsep{4pt}
\begin{tabular*}{\textwidth}{@{\extracolsep\fill}p{2.4cm}llp{2.2cm}p{2.2cm}p{2.0cm}}
\toprule
\textbf{System} & \textbf{Yr} & \textbf{Venue} &
\textbf{SAR Source} & \textbf{SAR Tasks} & \textbf{Limitation} \\
\midrule
EarthGPT~\cite{earthgpt2025}
  & 2025 & IEEE TGRS & MMRS-1M (multi)
  & VQA, Det & Below optical \\
SARChat-Bench-2M~\cite{sarchatchat2025}
  & 2025 & arXiv & Sentinel-1, GF-3
  & VQA, Class. & No OBB; eval only \\
SARLANG-1M~\cite{sarlang2026}
  & 2026 & TGRS & Multi-sensor SAR
  & VQA, Caption & No spatial output \\
SARCLIP~\cite{sarclip2026}
  & 2026 & ISPRS JP & SAR img-text
  & Retrieval & Single-modal only \\
SARVLM~\cite{sarvlm2026}
  & 2026 & [Pre] & SARLANG-1M & VQA, Det
  & Limited task range \\
Earth-OneVision~\cite{earthonevision2026}
  & 2026 & [Pre] & MMRS-34M (multi) & All tasks
  & SAR OBB absent \\
\botrule
\end{tabular*}
\footnotetext{EarthGPT achieves the highest published SAR task performance
among RS-MLLMs but does not match specialised SAR detectors on standard SAR
detection benchmarks (SSDD, HRSID, SAR-Ship).
No published RS-MLLM achieves OBB detection in SAR imagery~\cite{sarchatchat2025b}.}
\end{table}

Several structural observations emerge from Table~\ref{tab:sar_mllm2}. First,
no SAR-specific RS-MLLM achieves competitive performance on standard SAR object
detection benchmarks (SSDD, HRSID, SAR-Ship)---a direct consequence of the
scattering-physics barrier. Second, the domain adaptation strategies deployed
so far (joint optical--SAR training, SAR-specific fine-tuning) teach the model
\emph{what} SAR images look like without encoding \emph{why} they look that way.
Third, no SAR-MLLM produces OBB output, despite ship detection in SAR being the
canonical OBB task in the SAR domain.

The physics-constrained SAR--MSI fusion approach developed in our prior
work~\cite{rayleigh2026} demonstrates that incorporating Rayleigh-scattering
constraints into the data-level fusion pipeline substantially improves
cross-sensor feature alignment. Extending this principle to the RS-MLLM
training objective---through a physics-informed loss term that penalises
responses inconsistent with SAR backscatter physics---represents a concrete
direction for closing the SAR performance gap. We return to this direction
in the future research agenda (Sect.~\ref{sec9}).

%% ──────────────────────────────────────────────────────────────
\subsection{Hyperspectral MLLMs: The Final Frontier}\label{sec5:hsi}
%% ──────────────────────────────────────────────────────────────

Hyperspectral imaging (HSI) represents the most severe modality gap in the
RS-MLLM landscape. No dedicated HSI RS-MLLM has been published as of mid-2026,
despite HSI being a primary modality for mineral mapping, precision agriculture,
environmental monitoring, and military target discrimination.

\subsubsection{Technical Barriers for HSI-MLLMs}\label{sec5:hsibarriers}

\textbf{Dimensional incompatibility.}
An HSI data cube has dimensions $H\!\times\!W\!\times\!B$ where $B \in [200,\,400]$
spectral bands. Vision encoders expect $H\!\times\!W\!\times\!3$ (RGB) input.
The standard workaround---PCA dimensionality reduction from $B$ bands to 3
principal components---preserves variance but destroys spectral structure: the
three PCA components are linear mixtures of all $B$ bands and do not correspond
to any physically meaningful spectral interval~\cite{bioucas2013hyperspectral}.
The spectral information that makes HSI valuable is largely discarded in the
process.

\textbf{Spectral signature as primary discriminant.}
In optical RGB imagery, texture, shape, and colour provide the primary
discriminating features between object categories. In HSI, the spectral signature
of each pixel---its reflectance across 200--400 wavelengths---is the primary
discriminant~\cite{chen2014spectral}. Two surfaces that appear visually identical
in RGB can be uniquely distinguished by their HSI spectral signatures: for example,
healthy vegetation and drought-stressed vegetation have similar RGB appearances
but diverge significantly in their NIR and SWIR reflectance
profiles~\cite{yokoya2018hyperspectral}. An RS-MLLM with an RGB encoder discards
precisely the information needed to make this distinction.

Figure~\ref{fig:hsi_sig} illustrates spectral signatures for four land-cover
types across the 400--2500\,nm range, showing where RGB sampling is adequate
and where it fails critically.

\begin{figure}[h]
\centering
\begin{tikzpicture}[font=\small, >=Stealth, scale=0.84]
  %% Axes
  \draw[->,thick,gray!60] (0,0) -- (9.5,0)
    node[right,font=\footnotesize,gray!70] {Wavelength (nm)};
  \draw[->,thick,gray!60] (0,0) -- (0,5.5)
    node[above,font=\footnotesize,gray!70] {Reflectance (\%)};
  %% Y-axis ticks
  \foreach \y/\lbl in {0/0, 1/10, 2/20, 3/30, 4/40, 5/50}{
    \draw[gray!25] (0,\y) -- (9.3,\y);
    \node[left,font=\tiny,gray!60] at (-0.1,\y) {\lbl};
  }
  %% X-axis ticks
  \foreach \x/\lbl in {0.5/400, 2.0/700, 3.5/1000, 5.0/1400, 7.0/2000, 9.0/2500}{
    \draw[gray!25,thin] (\x,-0.05) -- (\x,5.2);
    \node[below,font=\tiny] at (\x,-0.12) {\lbl};
  }

  %% Vegetation signature — green with red-edge and NIR plateau
  \draw[green!60!black, very thick]
    (0.5,0.8)  .. controls (1.2,0.6)  .. (1.7,0.5)
               .. controls (2.0,0.4)  .. (2.2,0.6)
               .. controls (2.5,1.0)  .. (2.8,3.8)
               .. controls (3.5,4.2)  .. (4.5,4.0)
               .. controls (5.2,3.2)  .. (5.8,3.5)
               .. controls (6.5,2.5)  .. (7.0,2.8)
               .. controls (8.0,1.8)  .. (9.0,1.2);
  %% Bare soil — smoothly increasing
  \draw[orange!80!black, thick]
    (0.5,0.6)  .. controls (2.5,1.4)  .. (4.0,2.2)
               .. controls (5.5,2.8)  .. (7.0,3.2)
               .. controls (8.2,3.5)  .. (9.0,3.6);
  %% Water — high in blue, drops rapidly beyond 700nm
  \draw[blue!70, thick]
    (0.5,2.4)  .. controls (1.5,2.8)  .. (2.1,2.2)
               .. controls (2.5,1.2)  .. (3.0,0.4)
               .. controls (4.5,0.1)  .. (9.0,0.05);
  %% Urban / concrete — flat, slightly increasing
  \draw[gray!65, thick, dashed]
    (0.5,1.0)  .. controls (3.0,1.6)  .. (5.0,1.9)
               .. controls (7.0,2.1)  .. (9.0,2.2);

  %% RGB band shading (approximate)
  \fill[blue!20,  opacity=0.35] (0.5,0) rectangle (1.2,5.0);
  \fill[green!20, opacity=0.35] (1.5,0) rectangle (2.2,5.0);
  \fill[red!20,   opacity=0.35] (2.2,0) rectangle (2.8,5.0);
  %% RGB labels
  \node[blue!70,  font=\tiny, rotate=90] at (0.85,2.5) {B};
  \node[green!60!black, font=\tiny, rotate=90] at (1.85,2.5) {G};
  \node[red!70,   font=\tiny, rotate=90] at (2.5,2.5)  {R};
  %% RGB annotation
  \draw[red!70, very thick, decorate,
        decoration={brace,amplitude=4pt}]
    (0.5,-0.5) -- (2.8,-0.5);
  \node[red!70, font=\tiny\bfseries] at (1.65,-0.85)
    {RGB (3 bands only)};

  %% NIR annotation
  \draw[<-,green!50!black, thick] (3.1,4.1) -- (3.8,4.5)
    node[right,font=\tiny,green!60!black,align=left]
    {NIR plateau:\\\emph{healthy veg. only}};
  %% Red-edge annotation
  \draw[<-,green!50!black, thick] (2.5,2.2) -- (1.8,3.5)
    node[left,font=\tiny,green!60!black,align=right]
    {Red edge\\\emph{(invisible to RGB)}};

  %% Legend
  \draw[green!60!black, very thick] (6.5,5.2) -- (7.2,5.2);
  \node[right,font=\scriptsize,green!60!black] at (7.25,5.2) {Vegetation};
  \draw[orange!80!black, thick]     (6.5,4.8) -- (7.2,4.8);
  \node[right,font=\scriptsize,orange!80!black] at (7.25,4.8) {Bare soil};
  \draw[blue!70, thick]             (6.5,4.4) -- (7.2,4.4);
  \node[right,font=\scriptsize,blue!70] at (7.25,4.4) {Water};
  \draw[gray!65, thick, dashed]     (6.5,4.0) -- (7.2,4.0);
  \node[right,font=\scriptsize,gray!65] at (7.25,4.0) {Urban};
\end{tikzpicture}
\caption{Illustrative spectral signatures of four RS land-cover types across
400--2500\,nm, generated for conceptual comparison based on
well-established spectral properties documented in the
literature~\cite{bioucas2013hyperspectral,yokoya2018hyperspectral,chen2014spectral}.
Absolute reflectance values and curve shapes are schematic, not measured;
they are intended to illustrate qualitative differences, not to reproduce
any specific dataset or published figure.
The RGB sampling window (shaded columns) captures only three broad bands;
the vegetation NIR plateau ($\approx\!800$--1300\,nm) and the red edge
($\approx\!700$--750\,nm)---the two most diagnostically powerful features for
vegetation health assessment---are entirely outside the RGB range and
invisible to any RGB-trained encoder. An RS-MLLM using a CLIP encoder
therefore cannot distinguish healthy vegetation from drought-stressed
vegetation, or separate vegetation from bare soil under certain lighting
conditions---tasks that HSI sensors solve trivially}\label{fig:hsi_sig}
\end{figure}

\subsubsection{HM-Bench: Empirical Confirmation of the HSI Gap}\label{sec5:hmbench}

HM-Bench~\cite{hmbench2026}, released in April 2026, provides the first
controlled empirical evaluation of MLLM performance on HSI-specific tasks.
The benchmark comprises 19,337 expert-annotated samples from public HSI
datasets, organised into 13 task categories: spectral unmixing, material
identification, anomaly detection, multi-temporal change detection, spatial
classification, band selection, target detection, vegetation analysis, soil
mapping, water quality assessment, urban land-cover mapping, mineral mapping,
and spatial--spectral reasoning.

The HM-Bench evaluation of 18 representative MLLMs yields an unambiguous
finding: all 18 evaluated models show significant difficulty on complex
spatial--spectral reasoning tasks, including spectral unmixing and change
detection~\cite{hmbench2026}. This is not a marginal performance gap---it is
a systematic failure across all tested models and all model families
(CLIP-based, instruction-tuned, vision-language). The mechanistic explanation
is consistent with the dimensional incompatibility analysis above: models
apply RGB visual priors to PCA-compressed HSI channels, producing responses
driven by colour appearance rather than spectral physics. This failure mode
is the empirical foundation of the Type~III (Spectral) hallucination introduced
in Sect.~\ref{sec7}.

%% ──────────────────────────────────────────────────────────────
\subsection{Modality Distribution in RS-MLLM Literature}\label{sec5:dist}
%% ──────────────────────────────────────────────────────────────

Figure~\ref{fig:modality_dist} visualises the distribution of RS-MLLM papers
across sensor modalities, based on the quantitative synthesis of the 210 papers included
in this systematic review.

\begin{figure}[h]
\centering
\begin{tikzpicture}[font=\small, >=Stealth, scale=0.78]
  %% ── Left: Pie chart — overall distribution ─────────────────
  \def\pieR{2.0}
  %% Optical: 85% = 306 deg, SAR: 12% = 43.2 deg, HSI: 3% = 10.8 deg
  %% Draw pie slices
  %% Optical slice (306 deg = from 90 to -216 = 90 to 144 going anticlockwise)
  \fill[teal!65]
    (0,0) -- ({cos(90)*\pieR}, {sin(90)*\pieR})
    arc[start angle=90, end angle=90-306, radius=\pieR] -- cycle;
  %% SAR slice (43.2 deg after optical ends = from -216 to -259.2)
  \fill[orange!60]
    (0,0) -- ({cos(90-306)*\pieR}, {sin(90-306)*\pieR})
    arc[start angle=90-306, end angle=90-306-43.2, radius=\pieR] -- cycle;
  %% HSI slice (remainder = 10.8 deg)
  \fill[purple!55]
    (0,0) -- ({cos(90-349.2)*\pieR}, {sin(90-349.2)*\pieR})
    arc[start angle=90-349.2, end angle=90, radius=\pieR] -- cycle;
  \draw[white, line width=1.5pt] (0,0) circle (\pieR);
  %% Pie labels
  \node[white, font=\bfseries\small] at (-0.4,0.6) {Optical};
  \node[white, font=\bfseries\small] at (-0.4,0.0) {84.6\,\%};
  \node[font=\bfseries\scriptsize]   at (1.9,-1.7) {SAR 11.9\,\%};
  \node[purple!70,font=\bfseries\scriptsize] at (0.3,2.1) {HSI 2.9\,\%};
  \node[font=\scriptsize\bfseries]   at (0,-2.7) {(a) Overall distribution ($n\!=\!210$)};

  %% ── Right: Bar chart — papers per year per modality ────────
  \begin{scope}[xshift=6.0cm]
    %% Axes
    \draw[->,thick,gray!60] (0,0) -- (0,4.8)
      node[above,font=\tiny,gray!60] {\#\,Papers};
    \draw[->,thick,gray!60] (-0.2,0) -- (5.8,0)
      node[right,font=\tiny,gray!60] {Year};
    %% Y-axis ticks
    \foreach \y/\lbl in {0/0,1/5,2/10,3/20,4/35,4.6/50}{
      \draw[gray!25] (0,\y) -- (5.5,\y);
      \node[left,font=\tiny,gray!55] at (-0.1,\y) {\lbl};
    }
    %% X-axis labels
    \foreach \x/\lbl in {0.6/2022,1.5/2023,2.4/2024,3.3/2025,4.2/2026}{
      \node[below,font=\tiny] at (\x+0.3,-0.12) {\lbl};
    }
    %% Data (approximate from meta-analysis of 210 papers):
    %% Year: 2022 2023 2024 2025 2026(partial)
    %% Optical: 3   8   24  42   20   (scaled: /50*4.6)
    %% SAR:     0   1   4   12   8
    %% HSI:     0   0   0   3    3
    \def\bw{0.22}
    %% 2022
    \fill[teal!65]    (0.50,0) rectangle +(\bw, 0.276);
    \fill[orange!60]  (0.74,0) rectangle +(\bw, 0.000);
    \fill[purple!55]  (0.98,0) rectangle +(\bw, 0.000);
    %% 2023
    \fill[teal!65]    (1.40,0) rectangle +(\bw, 0.736);
    \fill[orange!60]  (1.64,0) rectangle +(\bw, 0.092);
    \fill[purple!55]  (1.88,0) rectangle +(\bw, 0.000);
    %% 2024
    \fill[teal!65]    (2.30,0) rectangle +(\bw, 2.208);
    \fill[orange!60]  (2.54,0) rectangle +(\bw, 0.368);
    \fill[purple!55]  (2.78,0) rectangle +(\bw, 0.000);
    %% 2025
    \fill[teal!65]    (3.20,0) rectangle +(\bw, 3.864);
    \fill[orange!60]  (3.44,0) rectangle +(\bw, 1.104);
    \fill[purple!55]  (3.68,0) rectangle +(\bw, 0.276);
    %% 2026 (partial year)
    \fill[teal!65]    (4.10,0) rectangle +(\bw, 1.840);
    \fill[orange!60]  (4.34,0) rectangle +(\bw, 0.736);
    \fill[purple!55]  (4.58,0) rectangle +(\bw, 0.276);
    %% Outlines
    \foreach \x in {0.50,1.40,2.30,3.20,4.10}{
      \draw[gray!40,thin]
        (\x,0) rectangle +(\bw+0.48, 4.8);
    }
    %% Legend
    \fill[teal!65]   (0.0,4.5) rectangle +(0.3,0.22);
    \node[right,font=\tiny] at (0.35,4.61) {Optical};
    \fill[orange!60] (0.0,4.15) rectangle +(0.3,0.22);
    \node[right,font=\tiny] at (0.35,4.26) {SAR};
    \fill[purple!55] (0.0,3.80) rectangle +(0.3,0.22);
    \node[right,font=\tiny] at (0.35,3.91) {HSI};
    \node[font=\scriptsize\bfseries] at (2.9,-0.5)
      {(b) Papers per year per modality};
  \end{scope}
\end{tikzpicture}
\caption{Modality distribution in the RS-MLLM literature ($n\!=\!210$ papers,
Jan 2022 -- Aug 2026). Panel~(a): Optical modality accounts for 84.6\,\% of
all RS-MLLM papers; SAR for 11.9\,\%; HSI for 2.9\,\%, based on the primary
sensor modality addressed. Panel~(b): Year-by-year breakdown shows accelerating
SAR activity from 2025 onward and nascent HSI activity beginning in 2025--2026.
Counts from the quantitative synthesis of 210 included papers
(Sect.~\ref{sec2:prisma}).}\label{fig:modality_dist}
\end{figure}

Figure~\ref{fig:modality_dist} confirms the optical dominance quantitatively.
Panel~(b) reveals two important temporal trends: first, SAR-focused RS-MLLM
papers accelerate sharply in 2025--2026, driven by SARChat-Bench-2M and SARLANG-1M;
second, HSI activity begins only in 2025--2026 and at a much lower absolute
rate, confirming that HSI remains the least-developed modality in the RS-MLLM
landscape. The absence of HSI papers prior to 2025 is consistent with the
HM-Bench finding: the evaluation infrastructure needed to demonstrate and
measure HSI failure modes did not exist until HM-Bench was released in April 2026.

%% ──────────────────────────────────────────────────────────────
\subsection{Cross-Modal Fusion Gaps}\label{sec5:fusion}
%% ──────────────────────────────────────────────────────────────

The most sophisticated RS applications require simultaneous reasoning across
multiple sensor modalities: SAR provides all-weather structural coverage;
optical provides spectral colour context; HSI provides material-composition
fingerprints. As established in the modality coverage analysis of
Sect.~\ref{sec3:modality}, the cross-modal fusion matrix has five of nine
cells entirely empty---all involving HSI.

\textbf{SAR--Optical fusion} is the most practically valuable cross-modal RS
task and the most developed. EarthGPT~\cite{earthgpt2025} is the first published
RS-MLLM to explicitly train on paired SAR--optical data within a unified
model (MMRS-1M), though its fusion strategy is simple modality-specific
feature concatenation without any physics-informed constraint.
Earth-OneVision~\cite{earthonevision2026} substantially extends multi-sensor
scope to six modalities (optical, SAR, infrared, multispectral, temporal, and
video), covering nine task categories---representing the most comprehensive
sensor and task coverage of any published RS-MLLM. However, Earth-OneVision
similarly uses a unified optical-domain encoder for all modalities without
physics-informed specialisation, meaning that SAR backscatter physics and
multispectral spectral properties are not explicitly encoded. The Rayleigh-constrained
approach in our prior work~\cite{rayleigh2026} demonstrates that incorporating SAR
backscatter physics into the data-level fusion pipeline significantly improves
cross-sensor alignment compared to this naive concatenation strategy.

\textbf{HSI--Optical fusion} has zero published RS-MLLM implementations.
The technical barriers are threefold: (i)~no large-scale co-registered
HSI--optical dataset exists at the scale required for MLLM training;
(ii)~no HSI-compatible encoder exists that can be paired with an optical
MLLM encoder; and (iii)~the Type~III spectral hallucination problem
(Sect.~\ref{sec7}) would corrupt HSI representations even if a fusion
architecture were available.

\textbf{SAR--HSI fusion} similarly has no published RS-MLLM implementations.
While SAR--HSI fusion has been studied for hyperspectral classification
using conventional deep learning~\cite{hong2021multimodal}, extending this
to the MLLM framework requires both a physics-aware SAR encoder and a
spectral-physics-aware HSI encoder---neither of which exists in published form.

\subsection*{Summary}

The multi-modal coverage analysis reveals a three-tier structure in RS-MLLM
development maturity. Optical RS-MLLMs are approaching performance saturation
on existing benchmarks and represent a well-established research direction.
SAR-MLLMs are nascent but active: five to six systems exist, driven by the
availability of SARChat-Bench-2M and SARLANG-1M, but none addresses the
scattering-physics barrier with a dedicated physics-informed encoder, and none
achieves OBB detection in SAR imagery.
Although Earth-OneVision~\cite{earthonevision2026} substantially expands sensor
and task coverage---processing optical, SAR, infrared, multispectral, temporal,
and video modalities across nine task categories---it uses a shared
optical-domain encoder for all modalities without physics-specific encoding,
and does not support HSI.
HSI-MLLMs do not yet exist as a dedicated model class: HM-Bench provides
evaluation infrastructure and empirical confirmation of systematic failure across
18 tested models, but no trained HSI-dedicated RS-MLLM has been published.
The remaining structural gaps---broad HSI coverage, physics-specific modality
encoding (SAR backscatter, spectral unmixing), systematic cross-sensor
evaluation protocols, and operationally grounded reasoning---cannot be closed
by scaling multi-sensor training data alone; they require modality-specific
encoder design and physics-informed training objectives, as analysed in the
future directions of Sect.~\ref{sec9}.
